% =================================================
% Справочный лист
% =================================================

\begin{center}
  {\fontsize{25}{30}\selectfont\bfseries\sffamily\color{Primary}
    Точка, прямая, отрезок и луч}

  \vspace{0.3em}

  {\large\sffamily\color{GrayText}
    Первые геометрические объекты и их обозначения}
\end{center}

\vspace{0.45em}

\section{Точка}

\begin{propertybox}[Определение]
  \textbf{Точка} указывает положение. На рисунке её изображают маленькой
  меткой и обозначают заглавной латинской буквой.
\end{propertybox}

\pointsexamplefigure

На рисунке отмечены точки \(A\), \(B\) и \(C\). Сама буква не является
точкой: она лишь служит её названием.

\section{Прямая}

\begin{propertybox}[Главное свойство]
  \textbf{Прямая} не имеет ни начала, ни конца. На рисунке можно показать
  только часть прямой.
\end{propertybox}

\lineexamplefigure

Прямую можно обозначить:
\begin{itemize}
  \item одной строчной латинской буквой: прямая \(a\);
  \item двумя точками, лежащими на ней: прямая \(AB\).
\end{itemize}

\begin{notebox}[Школьная запись]
  У прямой на рисунке не ставят стрелки. Линия просто продолжается за
  отмеченные точки в обе стороны.
\end{notebox}

\begin{propertybox}[Через две точки]
  Через любые две различные точки можно провести прямую, и притом только
  одну.
\end{propertybox}

\section{Точки на прямой и вне прямой}

\incidenceexamplefigure

По рисунку можно сказать:
\begin{itemize}
  \item точки \(A\) и \(B\) лежат на прямой \(a\);
  \item прямая \(a\) проходит через точки \(A\) и \(B\);
  \item точка \(C\) не лежит на прямой \(a\).
\end{itemize}

\section{Отрезок}

\begin{propertybox}[Определение]
  Точки \(A\) и \(B\) вместе со всеми точками прямой, расположенными между
  ними, образуют \textbf{отрезок} \(AB\). Точки \(A\) и \(B\) называют
  концами отрезка.
\end{propertybox}

\segmentexamplefigure

Отрезок имеет два конца, поэтому его можно измерить. Отрезки \(AB\) и
\(BA\) — один и тот же отрезок.

\begin{warningbox}[Не путайте объект и его длину]
  \(AB\) может обозначать отрезок, а запись \(AB=5\text{ см}\) сообщает
  длину этого отрезка.
\end{warningbox}

\section{Луч}

\begin{propertybox}[Определение]
  \textbf{Луч} имеет начало, но не имеет конца. Луч \(AB\) начинается в
  точке \(A\), проходит через точку \(B\) и продолжается дальше.
\end{propertybox}

\rayexamplefigure

Первая буква в названии луча всегда обозначает его начало. Поэтому лучи
\(AB\) и \(BA\) обычно различны.

\begin{notebox}[Как обозначаем луч]
  Пишем \textbf{луч \(AB\)} — без стрелки над буквами. Знак
  \(\overrightarrow{AB}\) в дальнейшем будет использоваться для вектора,
  а не для названия луча.
\end{notebox}

\oppositeraysfigure

Лучи \(AB\) и \(AC\) имеют общее начало \(A\) и идут в разные стороны.
Такие лучи называют противоположными.

\section{Сравним четыре объекта}

\objectscomparisonfigure

\begin{center}
  \renewcommand{\arraystretch}{1.38}
  \begin{tabularx}{\textwidth}{
    >{\raggedright\arraybackslash}p{0.20\textwidth}
    >{\centering\arraybackslash}p{0.20\textwidth}
    >{\raggedright\arraybackslash}p{0.32\textwidth}
    >{\centering\arraybackslash}X}
    \toprule
    \rowcolor{TableHeader}
    Объект & Сколько концов? & Как продолжается? & Можно измерить длину? \\
    \midrule
    Точка & — & — & нет \\
    Прямая & нет & в обе стороны & нет \\
    Отрезок & два & не продолжается за концы & да \\
    Луч & один — начало & в одну сторону & нет \\
    \bottomrule
  \end{tabularx}
\end{center}

\begin{warningbox}[Что определяет смысл записи \(AB\)]
  Одни и те же буквы могут встречаться в названиях разных объектов. Поэтому
  важно произносить полное название: \textit{прямая \(AB\)},
  \textit{отрезок \(AB\)} или \textit{луч \(AB\)}.
\end{warningbox}

\section{Проверьте себя}

\begin{questionsbox}
  \begin{enumerate}
    \item Чем прямая отличается от отрезка?
    \item Какая буква в названии луча указывает его начало?
    \item Почему отрезки \(AB\) и \(BA\) совпадают, а лучи \(AB\) и
      \(BA\) обычно не совпадают?
    \item Сколько прямых можно провести через две различные точки?
  \end{enumerate}
\end{questionsbox}
